\begin{helper}
Let $G$ be a finite $\Delta$-regular graph, let $r$ be a vertex, and let
$X\subseteq N_G(r)$.  Assume that whenever $u,v\in X$ are distinct, every
common neighbor of $u$ and $v$ outside $r$ is also adjacent to $r$.  Let
$\mathcal J_2$ be the independent two-element subsets of $X$ and let
$\mathcal J_3$ be the independent three-element subsets of $X$.  For a
two-element set $P\in\mathcal J_2$, let
\[
C(P)=\{w:\text{$w$ is adjacent in $G$ to every vertex of $P$}\}.
\]
If $W(P)\ge0$ for every $P\in\mathcal J_2$, then
\[
\sum_{Q\in\mathcal J_3}\ \sum_{\substack{P\subseteq Q\\ |P|=2}} W(P)
\le
\sum_{P\in\mathcal J_2}\bigl(\Delta-|C(P)|\bigr)W(P).
\]
\end{helper}
\begin{proof}
First fix $P\in\mathcal J_2$, and write $P=\{u,v\}$ with $u\ne v$.  Because
$P\subseteq X\subseteq N_G(r)$, the root $r$ is adjacent to both $u$ and
$v$, so $r\in C(P)$.  If $w\in C(P)$ and $w\ne r$, then $w$ is a common
neighbor of the distinct vertices $u,v\in X$.  The clean-neighborhood
hypothesis therefore implies that $w$ is adjacent to $r$.  Also $w\notin P$,
since otherwise $w$ would be adjacent to itself.  Hence
\[
C(P)\setminus\{r\}\subseteq N_G(r)\setminus P. \tag{1}
\]

Let $\mathcal F(P)=\{Q\in\mathcal J_3:P\subseteq Q\}$.  For every
$Q\in\mathcal F(P)$, the set $Q\setminus P$ has size one, because
$|Q|=3$, $|P|=2$, and $P\subseteq Q$.  Let $w_Q$ be its unique element.
Then $Q=P\cup\{w_Q\}$.  Since $Q\subseteq X$, we have $w_Q\in X$; since
$w_Q\notin P$ by construction; and since $Q$ is independent, $w_Q$ is not
adjacent to $u$, so $w_Q\notin C(P)$.  Thus
\[
w_Q\in X\setminus(C(P)\cup P).
\]
Moreover the map $Q\mapsto w_Q$ is injective, because
$Q=P\cup\{w_Q\}$.  Therefore
\[
|\mathcal F(P)|\le |X\setminus(C(P)\cup P)|. \tag{2}
\]

The sets $X\setminus(C(P)\cup P)$ and $C(P)\setminus\{r\}$ are disjoint,
and by (1) and $X\subseteq N_G(r)$ their union is contained in
$N_G(r)\setminus P$.  Since $G$ is $\Delta$-regular at $r$ and $P$ is a
two-element subset of $N_G(r)$,
\[
|N_G(r)\setminus P|=\Delta-2.
\]
Also $r\in C(P)$, so $|C(P)\setminus\{r\}|=|C(P)|-1$.  Combining these
cardinality facts gives
\[
|X\setminus(C(P)\cup P)|+|C(P)\setminus\{r\}|
\le \Delta-2,
\]
and hence
\[
|X\setminus(C(P)\cup P)|\le \Delta-|C(P)|.
\]
Together with (2), this proves
\[
|\mathcal F(P)|\le \Delta-|C(P)|
\quad\text{for every }P\in\mathcal J_2. \tag{3}
\]

Now reindex the double sum by its distinguished pair:
\[
\sum_{Q\in\mathcal J_3}\sum_{\substack{P\subseteq Q\\ |P|=2}} W(P)
=
\sum_P |\mathcal F(P)|W(P).
\]
If a two-element $P$ occurs inside some $Q\in\mathcal J_3$, then
$P\subseteq X$ and $P$ is independent, so $P\in\mathcal J_2$; all other
pairs contribute zero.  For $P\in\mathcal J_2$, multiply (3) by the
nonnegative number $W(P)$ and sum over $P\in\mathcal J_2$.  This gives the
claimed inequality.
\end{proof}
